\chapter{Cadre général du projet}
\label{chap:cadre_general}

\section{Introduction}
Ce premier chapitre a pour vocation de poser les fondements contextuels et organisationnels dans lesquels s'inscrit le projet de fin d'études. Avant de plonger dans les détails techniques et architecturaux de l'application, il est primordial de comprendre l'environnement métier, les contraintes industrielles et les défis quotidiens auxquels le Groupe COFAT fait face. Nous commencerons par une présentation exhaustive de l'entreprise, de son histoire depuis 1984, de ses produits phares et de son envergure internationale couvrant six pays et dix entités distinctes. Par la suite, nous réaliserons une étude critique de l'existant, en décortiquant les processus actuels de gestion des capacités (équipements, espaces et ressources humaines) afin d'en extraire les limites et les failles. Cette analyse nous conduira naturellement à formuler la problématique centrale du projet. En réponse à ces enjeux, nous dévoilerons les contours de la solution numérique proposée, \textit{Cofat Capacity Study}, avant de détailler la méthodologie de gestion de projet agile (Scrum) adoptée pour garantir la réussite de cette réalisation.

\section{Présentation de l'organisme d'accueil : Le Groupe COFAT}

\subsection{Historique et cœur de métier}
Fondé en 1984, le Groupe COFAT (Compagnie Faisceaux Autos Tunisiens) est l'une des filiales majeures du Groupe Elloumi, un conglomérat industriel de renommée internationale. Depuis sa création, COFAT s'est spécialisé dans un domaine d'ingénierie hautement technique et exigeant : la conception, le développement et la fabrication de faisceaux de câbles pour l'industrie automobile et les véhicules utilitaires. 

Le faisceau de câbles, souvent qualifié de « système nerveux » du véhicule, est un assemblage complexe de fils électriques, de connecteurs, de gaines de protection, de terminaux et d'agrafes. Son rôle est de distribuer l'énergie électrique et de transmettre les signaux d'information entre l'ensemble des composants électroniques et mécaniques du véhicule (moteur, capteurs, éclairage, systèmes d'infodivertissement, calculateurs, etc.). La fabrication de ces produits requiert une précision millimétrique et obéit à des normes de sécurité drastiques, car une défaillance d'un faisceau peut entraîner la panne complète du véhicule ou compromettre la sécurité des passagers.

Au fil des décennies, COFAT a su diversifier ses segments de produits pour répondre aux évolutions technologiques de l'automobile. Aujourd'hui, l'entreprise produit :
\begin{itemize}
    \item \textbf{Les faisceaux habitacles} : Ils relient les éléments internes de la cabine (tableau de bord, sièges, portes, climatisation) et gèrent le confort et la sécurité passive.
    \item \textbf{Les faisceaux moteur} : Conçus pour résister à des environnements extrêmes (hautes températures, vibrations intenses, projections de fluides corrosifs), ils assurent le fonctionnement de la motorisation et des systèmes de freinage (ABS, ESP).
    \item \textbf{Les câbles haute tension} : Avec l'avènement des véhicules électriques (BEV) et hybrides (PHEV), COFAT a investi dans la R\&D pour produire des câblages capables de supporter des tensions et des courants très élevés, indispensables à la propulsion électrique.
\end{itemize}

Grâce à son expertise technique, sa flexibilité industrielle et sa capacité d'innovation, COFAT a su gagner la confiance des plus grands constructeurs automobiles mondiaux. Ses clients principaux comptent aujourd'hui des marques prestigieuses telles que Stellantis, Volkswagen, Renault, Scania, Volvo, Daimler et PSA.

\subsection{Envergure internationale et entités du Groupe}
Pour accompagner ses clients au plus près de leurs sites d'assemblage et pour optimiser sa chaîne logistique, COFAT a déployé une stratégie d'internationalisation ambitieuse. Aujourd'hui, le groupe compte plus de 5 600 collaborateurs répartis sur 10 entités opérationnelles situées dans 6 pays différents. Cette présence mondiale se structure autour de différents types de sites, chacun ayant une vocation stratégique spécifique :

\begin{itemize}
    \item \textbf{PLANT (Sites de production)} : Ces usines sont le cœur industriel du groupe. Elles abritent les lignes de coupe, de préparation et d'assemblage des faisceaux de câbles. Elles sont dimensionnées pour produire de grands volumes avec une qualité zéro défaut.
    \item \textbf{D\&D (Centres de Développement \& Design)} : Ces centres regroupent les ingénieurs d'études. Ils travaillent en amont des projets avec les constructeurs pour concevoir les schémas électriques, optimiser le routage des câbles dans le véhicule et prototyper les nouveaux produits.
    \item \textbf{Backoffice et IT} : Ces sites centralisent l'administration, les finances, les ressources humaines au niveau corporate, ainsi que la gestion de l'infrastructure informatique et des systèmes d'information du groupe.
    \item \textbf{SERVICE (Centres de support)} : Situés à proximité stratégique des usines des clients finaux, ils assurent l'interface commerciale, la logistique \textit{Just-In-Time} et le support technique immédiat.
\end{itemize}

Le tableau \ref{tab:cofat_sites} récapitule l'ensemble des implantations mondiales du Groupe COFAT.

\begin{table}[h]
\centering
\begin{tabular}{|l|l|l|l|}
\hline
\textbf{Entité} & \textbf{Pays} & \textbf{Ville} & \textbf{Type de site} \\ \hline
Cofat Tunis (Backoffice) & Tunisie & Kram, Tunis & Backoffice \\ \hline
Cofat Tunis (D\&D) & Tunisie & Sidi Hssin, Tunis & D\&D \\ \hline
Cofat Mateur & Tunisie & Mateur & PLANT \\ \hline
Cofat Kairouan & Tunisie & Kairouan & PLANT \\ \hline
ABM Informatique & Tunisie & Tunis & IT Services \\ \hline
Cofat Deutschland & Allemagne & Gifhorn & D\&D \\ \hline
Cofat Egypt & Égypte & Le Caire & PLANT \\ \hline
Cofat Mexico & Mexique & Durango & D\&D \\ \hline
Cofat Brasil & Brésil & São Paulo & PLANT \\ \hline
Cofat USA & États-Unis & El Paso, Texas & SERVICE \\ \hline
\end{tabular}
\caption{Répartition mondiale des entités du Groupe COFAT}
\label{tab:cofat_sites}
\end{table}

\subsection{Management de la qualité et certifications}
Dans l'industrie automobile, la qualité n'est pas une option, c'est une condition de survie. Chaque composant doit garantir une fiabilité absolue. Pour structurer ses processus, homogénéiser ses méthodes sur l'ensemble de ses 10 sites et rassurer ses clients donneurs d'ordres, COFAT s'appuie sur un Système de Management Intégré (SMI) rigoureux, validé par plusieurs certifications internationales majeures :

\begin{itemize}
    \item \textbf{IATF 16949} : Il s'agit du standard de management de la qualité le plus exigeant et spécifique à l'industrie automobile. Il met l'accent sur la prévention des défauts, la réduction des variations et la diminution des gaspillages dans la chaîne d'approvisionnement.
    \item \textbf{ISO 9001} : C'est la norme de base du système de management de la qualité globale, garantissant que l'entreprise fournit de manière constante des produits et services conformes aux exigences des clients et de la réglementation.
    \item \textbf{ISO 14001} : Cette certification atteste de l'engagement de COFAT en matière de management environnemental, en limitant l'impact de ses activités industrielles, en optimisant la gestion de ses déchets (notamment les chutes de câbles en cuivre) et en maîtrisant sa consommation énergétique.
    \item \textbf{ISO 45001} : Ce standard international certifie le système de management de la santé et de la sécurité au travail. Il démontre la volonté de la direction d'offrir à ses 5 600 employés un environnement de travail sûr, en minimisant les risques d'accidents (coupures, troubles musculosquelettiques liés à l'assemblage manuel).
\end{itemize}

\section{Étude de l'existant}
Avant le lancement du projet \textit{Cofat Capacity Study}, la gestion des capacités de production au sein du Groupe COFAT reposait sur des méthodes traditionnelles, fragmentées et chronophages. Bien que le groupe dispose d'outils ERP pour la gestion de la production courante, la planification stratégique à moyen et long terme — c'est-à-dire l'anticipation des besoins en machines, en surfaces au sol et en main-d'œuvre pour les nouveaux projets automobiles (comme le lancement d'une nouvelle série VW ou Stellantis) — était gérée de manière totalement décentralisée.

Le processus s'articulait autour de trois volets d'analyse majeurs, gérés indépendamment sur chaque site industriel (Mateur, Kairouan, São Paulo, Le Caire) à l'aide de dizaines de fichiers tableurs Microsoft Excel.

\subsection{Le volet Équipements}
Pour qu'une usine puisse assembler les faisceaux d'un nouveau projet, elle doit s'assurer de disposer des équipements industriels nécessaires, allant des simples presses de sertissage aux complexes machines de coupe automatique (Komax), en passant par les tables de test électrique et les machines de soudure ultrason.

Le fichier Excel utilisé localement par chaque \textit{Site Manager} comportait plusieurs colonnes critiques :
\begin{itemize}
    \item \textbf{machineNeed} : Le besoin théorique en machines, calculé en fonction du temps de cycle par faisceau, des volumes prévisionnels commandés par le client et du nombre d'équipes (shifts) prévues.
    \item \textbf{availableMachine} : Le parc matériel physiquement disponible et fonctionnel sur le site, hors machines en maintenance curative lourde.
    \item \textbf{load (Taux de charge)} : Un indicateur clé calculé par la formule $Load = (machineNeed / availableMachine) \times 100$. Un taux supérieur à 100\% indique une saturation et une incapacité à produire la demande.
    \item \textbf{toOrder} : La différence mathématique générant le nombre de machines supplémentaires à acquérir pour combler le déficit de capacité.
\end{itemize}
Bien que cette logique de calcul soit robuste, le fait qu'elle soit isolée dans des fichiers Excel non partagés empêchait le siège (Backoffice) de savoir si une machine manquante à Mateur pouvait être transférée depuis Kairouan, où elle serait potentiellement en sous-charge.

\subsection{Le volet Espaces}
L'industrie du câblage automobile est très gourmande en espace physique (m²). Les tables d'assemblage (jig boards) peuvent mesurer jusqu'à plusieurs mètres de long et nécessitent des allées larges pour l'approvisionnement logistique en bord de ligne. 

Les fichiers Excel de gestion d'espace suivaient les zones de production réelles de l'usine :
\begin{itemize}
    \item \textbf{Zone de Coupe} : Où les bobines de fils sont découpées et dénudées.
    \item \textbf{Zone de Préparation} : Où s'effectuent les épissures, le sertissage manuel et le twistage.
    \item \textbf{Zones Projets Assemblage} : Des halls dédiés à des projets clients spécifiques (ex: \textit{VW Tayron}, \textit{International E44}, \textit{Stellantis J4U}, \textit{VW Golf 8}, \textit{LS e-mobility BEV-A}).
\end{itemize}
Pour chaque zone, l'ingénieur méthodes saisissait la \textbf{surface totale disponible} et calculait la \textbf{surface occupée} par les lignes existantes pour en déduire le \textbf{taux d'occupation}. La difficulté résidait dans l'incapacité à avoir une vue 3D consolidée du site et à modéliser rapidement l'impact d'un nouveau projet sur l'encombrement global.

\subsection{Le volet Ressources Humaines (RH)}
La confection de faisceaux de câbles reste une activité nécessitant beaucoup de main-d'œuvre manuelle pour l'enrubannage et l'insertion des terminaux dans les boîtiers. L'anticipation des recrutements est donc vitale.

Les tableaux Excel RH étaient structurés selon les catégories réelles de personnel :
\begin{itemize}
    \item \textbf{Direct} : Les opérateurs de production travaillant sur les machines de coupe et préparation.
    \item \textbf{Assemblage Direct} : Les opérateurs affectés spécifiquement à l'assemblage manuel sur les tables des projets clients.
    \item \textbf{Indirect} : Le personnel d'encadrement, de logistique, de méthodes, de qualité et de maintenance.
\end{itemize}
Ces fichiers projetaient les effectifs nécessaires mois par mois pour l'année en cours (MO 01 à MO 12 de 2025) et par trimestre pour les années futures (Q01 à Q04 pour 2026 et 2027), en fonction de la rampe de lancement (\textit{ramp-up}) des nouveaux véhicules.

\subsection{Le processus manuel de consolidation}
Le problème majeur de cet existant ne résidait pas dans la donnée elle-même, mais dans la manière dont elle circulait. Trimestriellement, ou lors de l'étude d'un projet majeur, le Backoffice de Tunis devait consolider la capacité mondiale. 
Le processus était le suivant : un email était envoyé aux responsables des différents sites (Tunisie, Égypte, Brésil, etc.) demandant l'envoi de leurs fichiers Excel de capacité mis à jour. Le siège recevait alors une multitude de fichiers, souvent avec des formats altérés, des colonnes renommées, des formules cassées ou des macros incompatibles. 
Un ingénieur central devait ensuite passer plusieurs jours à ouvrir ces fichiers un par un, nettoyer les données, aligner les formats, et faire des copiés-collés manuels vers une immense "feuille maître" (Master Sheet) afin d'agréger les chiffres mondiaux.

\section{Problématique}
L'analyse de l'existant décrite précédemment met en évidence les limites sévères d'un système de gestion de capacité basé sur l'échange de fichiers bureautiques dans un contexte industriel multinational. Cette méthode d'un autre temps engendre des dysfonctionnements profonds qui impactent la compétitivité et la rentabilité du Groupe COFAT. Nous pouvons synthétiser la problématique autour de quatre axes majeurs :

\subsection{Les erreurs de saisie et de manipulation}
L'intervention manuelle lors de la consolidation (copié-collé entre fichiers) multiplie les risques d'erreurs humaines. Une erreur de virgule dans le prix d'un équipement ou une formule écrasée dans le calcul du besoin machine peut entraîner l'approbation d'un investissement de plusieurs dizaines de milliers d'euros totalement injustifié, ou à l'inverse, bloquer le lancement d'une production par manque de matériel. Le format Excel ne permet pas de contraindre strictement le typage des données ou d'appliquer des validations croisées dynamiques (ex: s'assurer que le département saisi est bien parmi "IT, HR, QUALITY, PRODUCTION").

\subsection{La perte de traçabilité et de suivi d'historique}
Lorsqu'une demande de budget ou de capacité est envoyée par mail sous forme de fichier joint, elle perd immédiatement son statut de "donnée vivante". Il devient impossible de savoir qui a modifié quelle cellule, à quelle date, et pour quelle raison. Par exemple, lors des négociations budgétaires, si le département des achats modifie le coût unitaire d'une machine après avoir trouvé un nouveau fournisseur, le demandeur initial n'est pas notifié. Cette déconnexion crée des litiges internes, des malentendus, et une perte de l'historique décisionnel indispensable lors des audits de qualité (IATF 16949).

\subsection{Le temps de consolidation excessif}
Dans l'industrie automobile, la réactivité (Time-to-Market) est un avantage concurrentiel décisif. Lorsqu'un constructeur demande à COFAT si l'entreprise est capable d'absorber une hausse de production inattendue de 30\% sur un faisceau, la réponse doit être rapide. Avec l'ancien processus, la collecte des fichiers des 10 entités, le nettoyage des données et l'agrégation manuelle prenaient des jours, voire des semaines. Ce délai de latence retardait la validation des budgets et menaçait l'obtention de nouveaux marchés.

\subsection{Le manque de sécurité et de confidentialité}
Les fichiers Excel circulant par email ou stockés sur des répertoires partagés locaux ne bénéficient pas d'une gestion fine des droits d'accès. Des données extrêmement sensibles (coûts d'investissement, stratégies industrielles, prévisions d'embauche de centaines d'opérateurs) peuvent être lues, copiées ou altérées par des personnes non autorisées. Il manque cruellement un système de contrôle d'accès basé sur les rôles (RBAC) pour isoler les données entre les usines, et distinguer les droits de lecture, de création et de validation.

\begin{table}[h]
\centering
\begin{tabular}{|p{4cm}|p{5cm}|p{6cm}|}
\hline
\textbf{Axe problématique} & \textbf{Situation existante (Excel/Email)} & \textbf{Solution logicielle proposée} \\ \hline
\textbf{Fiabilité des données} & Risque élevé d'erreurs de copié-collé, formules cassées & Validation stricte via API, base de données relationnelle (SQL Server) \\ \hline
\textbf{Traçabilité} & Fichiers figés, pas d'historique des modifications & Sauvegarde horodatée, système de notifications, traçabilité des actions \\ \hline
\textbf{Temps de consolidation} & Plusieurs jours (collecte et agrégation manuelles) & Agrégation automatique en temps réel par le serveur Node.js \\ \hline
\textbf{Sécurité} & Données en clair dans des fichiers circulant par mail & Authentification JWT, hachage bcrypt, contrôle d'accès par rôle (Admin, Site Manager, Achat) \\ \hline
\end{tabular}
\caption{Comparatif : Existant vs Solution proposée}
\label{tab:comparatif_existant}
\end{table}

\section{La solution proposée : \textit{Cofat Capacity Study}}
Pour pallier ces défaillances et doter le Groupe COFAT d'un outil d'aide à la décision stratégique à la hauteur de ses ambitions, j'ai proposé de concevoir et développer de zéro une plateforme web applicative centralisée, baptisée \textbf{« Cofat Capacity Study »}. Cette solution moderne vise à transformer un processus bureaucratique lourd en un flux d'informations fluide, sécurisé et en temps réel. 

La solution s'articule autour de cinq grands composants fonctionnels :

\subsection{Portail web multi-utilisateurs et sécurisé}
L'application fournit un point d'accès unique via un navigateur web. Elle intègre un système d'authentification robuste (basé sur des jetons JWT sécurisés) qui identifie l'utilisateur, vérifie son site d'affectation et surtout son rôle métier (Administrateur, Site Manager, ou Achat). L'interface utilisateur est dynamique et adapte les menus et les autorisations d'action en fonction de ce rôle, garantissant que chaque employé ne voit et ne modifie que les données qui le concernent.

\subsection{Module d'import Excel automatisé}
Conscient que les ingénieurs locaux sont habitués à travailler sur des tableurs et que la saisie manuelle de centaines de lignes de données RH ou Machines dans un formulaire web serait rejetée, le système intègre un parseur Excel intelligent. Le Site Manager peut continuer à utiliser son fichier de travail habituel, mais il le "glisse-dépose" sur la plateforme. Le serveur Node.js (via la bibliothèque XLSX) va alors extraire, valider (types, colonnes obligatoires) et insérer ces données de manière structurée et massive (bulk insert) directement dans la base SQL Server en quelques secondes.

\subsection{Moteur de consolidation en temps réel}
C'est le cœur analytique de l'application. Fini le rassemblement manuel des données : dès qu'un site (par exemple, Durango au Mexique) met à jour son planning de machines, l'administrateur au siège à Tunis peut ouvrir son tableau de bord "Cofat Group" et voir les graphiques de consolidation mondiaux (effectifs totaux, surfaces occupées globales) se mettre à jour instantanément grâce aux requêtes d'agrégation SQL pré-optimisées.

\subsection{Référentiel des standards (StandardInvestment V2)}
Pour éviter que chaque usine estime ses coûts d'investissement avec des fournisseurs différents, l'application centralise un catalogue de référence. Il contient la liste exacte des machines autorisées (avec leur code, technologie, fournisseur homologué, capacité de production et coût unitaire validé par le département des achats). Ce référentiel permet de calculer automatiquement et de manière standardisée l'enveloppe budgétaire nécessaire pour combler le manque de machines (les \textit{toOrder}) calculé lors de la planification locale.

\subsection{Workflow du budget non-industriel}
Au-delà de la capacité de production directe, l'application gère les demandes d'investissements annexes (IT, bâtiment, qualité, logistique). Un module spécifique permet de digitaliser le cycle de vie de la demande budgétaire, facilitant l'interaction entre le demandeur (Site Manager) et le département des achats. Un système de notifications en temps réel vient fluidifier ce processus de négociation budgétaire, avec des possibilités d'export vers Excel ou PDF pour les réunions de validation de la direction générale.

\section{Méthodologie adoptée : L'approche Scrum}
La construction d'une application de cette envergure, touchant à des processus métier complexes et s'adressant à plusieurs catégories d'utilisateurs à l'international, ne peut se satisfaire d'une méthode de développement traditionnelle "en cascade" (Waterfall). En effet, les besoins des utilisateurs, notamment sur les formules de calcul de capacité et les vues de consolidation, sont susceptibles d'évoluer au fur et à mesure qu'ils découvrent le potentiel de l'outil numérique.

C'est pourquoi nous avons opté pour le framework \textbf{Scrum}, la méthodologie agile la plus répandue. Scrum permet de livrer de la valeur métier très tôt et de manière itérative, en construisant le produit par incréments successifs, appelés \textbf{Sprints}. 

\subsection{Les rôles Scrum du projet}
L'organisation de l'équipe projet a été structurée selon les trois piliers du framework Scrum, adaptés au contexte académique et industriel de ce PFE :
\begin{itemize}
    \item \textbf{Le Product Owner (PO)} : Incarné par mon encadrant industriel au sein de COFAT (M. Mohamed Said BOUCHOUICHA), il représente la "voix du client". Il a la responsabilité de définir la vision du produit, de recenser les exigences métier, et de maintenir et prioriser le Product Backlog pour maximiser la valeur livrée.
    \item \textbf{Le Scrum Master} : Rôle assuré conjointement avec le soutien de mon encadrant académique (M. Mohamed ALOUANI), garantissant le respect de la méthode, facilitant les cérémonies et aidant à lever les obstacles techniques ou organisationnels.
    \item \textbf{La Development Team} : Je (M. ZAKRAOUI) constituais l'équipe de développement, prenant en charge l'architecture logicielle, le développement frontend et backend, l'intégration continue, les tests et la documentation, transformant ainsi les items du backlog en fonctionnalités opérationnelles.
\end{itemize}

\subsection{Les cérémonies Scrum appliquées}
Tout au long de ces six mois de projet, notre rythme de travail a été ponctué par les événements Scrum, garantissant la transparence, l'inspection et l'adaptation :
\begin{itemize}
    \item \textbf{Le Sprint Planning} : Au début de chaque itération, nous sélectionnions les User Stories prioritaires du Product Backlog pour définir le Sprint Backlog, en s'assurant que l'objectif du sprint était clair, techniquement faisable et apportait une valeur directe (ex: "Permettre l'import d'un fichier Excel RH").
    \item \textbf{Le Daily Scrum} : Des points de synchronisation fréquents (réunions courtes) avec le Product Owner permettaient de faire le point sur l'avancement, de montrer les premiers écrans développés et de valider rapidement des choix ergonomiques ou algorithmiques.
    \item \textbf{La Sprint Review} : À la fin de chaque sprint, l'incrément logiciel (une version fonctionnelle de l'application) était présenté aux parties prenantes (ingénieurs méthodes, département achat) pour recueillir leurs retours (feedbacks). Ces retours étaient immédiatement convertis en améliorations pour le sprint suivant.
    \item \textbf{La Sprint Retrospective} : Une analyse interne sur nos méthodes de travail (ex: améliorer la couverture de tests sur les requêtes SQL complexes, optimiser l'utilisation de Git Flow pour la gestion des versions du code).
\end{itemize}

\subsection{Planification et découpage en Sprints}
Le projet, s'étalant sur une durée globale de six mois (de début mars à début septembre 2026), a été structuré autour de quatre Sprints majeurs, précédés d'un Sprint 0 d'initialisation :

\begin{itemize}
    \item \textbf{Sprint 0 (Mars 2026)} : Définition de l'architecture logicielle, choix technologiques, mise en place des environnements de développement et de la base de données relationnelle.
    \item \textbf{Sprint 1 (Avril 2026)} : Développement du socle technique fondamental (Authentification JWT sécurisée) et du tableau de bord d'administration (gestion des utilisateurs, gestion des sites mondiaux). Ce sprint fera l'objet de la seconde partie du Chapitre 3.
    \item \textbf{Sprint 2 (Mai - Juin 2026)} : Réalisation des modules de "Planification Locale". Il s'agit du développement des interfaces et des APIs permettant aux Site Managers d'importer et d'analyser leurs données d'Équipements, d'Espaces et de Ressources Humaines. 
    \item \textbf{Sprint 3 (Juillet 2026)} : Réalisation du module de "Consolidation Groupe" (agrégation des données des 10 sites en temps réel) et refonte du catalogue des standards industriels (StandardInvestment V2). Les Sprints 2 et 3 seront détaillés dans le Chapitre 4.
    \item \textbf{Sprint 4 (Août 2026)} : Conception et implémentation du module de "Budget Non-Industriel", intégrant le workflow collaboratif, les notifications temps réel et les exports de reporting (Excel/PDF). Ce dernier sprint sera couvert par le Chapitre 5.
\end{itemize}

\section{Conclusion}
À l'issue de ce premier chapitre, nous avons cerné avec précision le contexte industriel exigeant du Groupe COFAT et les enjeux liés à la gestion de la capacité de production de ses dix usines. L'étude de l'existant a révélé les failles d'un processus manuel basé sur des tableurs locaux, induisant lenteur, opacité et risques d'erreurs considérables. La nécessité d'une rupture technologique est flagrante. La plateforme \textit{Cofat Capacity Study} est la réponse directe à cette problématique. En s'appuyant sur la méthode agile Scrum, le développement de cette solution centralisée s'annonce itératif et orienté vers la création de valeur immédiate pour les utilisateurs. 
Le chapitre suivant se concentrera sur la phase d'analyse et de spécification, où nous traduirons ces défis organisationnels en exigences fonctionnelles claires, en identifiant les acteurs clés du système et en modélisant leurs interactions futures avec l'application.
